\chapter{Mapa de pesos e prioridades}

\section{Estrutura oficial da prova}

\begin{center}
\begin{tabular}{@{}p{2.2cm}p{5.2cm}rrr@{}}
\toprule
\textbf{Módulo} & \textbf{Disciplina} & \textbf{Questões} & \textbf{Peso} & \textbf{Pontos} \\
\midrule
I & Língua Portuguesa & 12 & 1 & 12 \\
I & Língua Inglesa & 12 & 1 & 12 \\
I & Raciocínio Lógico Matemático & 5 & 1 & 5 \\
I & Atualidades e Inteligência Artificial & 6 & 1 & 6 \\
I & Legislação (Seg.\ Informação e Proteção de Dados) & 5 & 1 & 5 \\
II & Conhecimentos Específicos & 30 & 2,5 & 75 \\
\midrule
& \textbf{Total} & \textbf{70} & & \textbf{115} \\
\bottomrule
\end{tabular}
\end{center}

\section{Distribuição real do Módulo II na Dataprev 2024 (30 questões)}

\begin{center}
\begin{tabular}{@{}p{6.4cm}rr@{}}
\toprule
\textbf{Bloco} & \textbf{Questões} & \textbf{Pontos} \\
\midrule
Engenharia de Software & 9 & 22,5 \\
Banco de Dados / BI & 6 & 15 \\
Programação & 6 & 15 \\
Arquitetura de Software & 4 & 10 \\
Segurança da Informação & 3 & 7,5 \\
Redes de Computadores & 2 & 5 \\
\bottomrule
\end{tabular}
\end{center}

\textit{Reclassificado questão a questão direto no PDF em 29/07/2026 --- a
versão anterior tinha categorias que nem existem como tag no quiz e não
fechava as 30 questões.}

\section{Comportamento da FGV em outras provas de TI}

\begin{center}
\begin{tabular}{@{}p{3.2cm}p{2cm}p{5.6cm}@{}}
\toprule
\textbf{Prova} & \textbf{Data} & \textbf{Blocos dominantes} \\
\midrule
MPU, Desenv.\ de Sistemas & mai/2025 & Português 15, \textbf{Banco de Dados 12}, Eng.\ Software 8 \\
AgSUS, Analista de TI & out/2025 & Português 12, Informática 8, RL 7 \\
TJ-RJ, Analista de Sistemas & jan/2026 & Português 19, Eng.\ Software 7, \textbf{Banco de Dados 7} \\
\bottomrule
\end{tabular}
\end{center}

\section{As cinco conclusões que moldam a priorização}

\begin{enumerate}
\item \textbf{Engenharia de Software, Banco de Dados/BI e Programação dividem
o topo.} Na Dataprev 2024 vieram quase empatados (9, 6 e 6 questões). No MPU
o BD disparou na frente de Engenharia de Software. Trate os três como
prioridade, sem eleger vencedor único --- e sem esquecer Arquitetura de
Software (4 questões, tão relevante quanto Segurança).
\item \textbf{Redes de Computadores não é tão ``fora do edital'' quanto
parece.} Das 2 questões de 2024, só a de X.800/arquitetura de segurança OSI
foge de fato do edital do perfil. A outra (ambientes de Internet, intranet,
extranet e portal) é o item 4 do próprio edital de Desenvolvimento de
Sistemas --- só não tem uma disciplina ``Redes'' com nome próprio. Ainda
vale treinar OSI/X.800 à parte: é o ponto realmente cego.
\item \textbf{Raciocínio Lógico da FGV é matemática, não lógica formal (mas
o edital 2026 pede as duas).} O que caiu em 2024 foi aritmética, álgebra,
estatística, análise combinatória e geometria plana. Só uma questão de
proposições. O edital 2026 lista literalmente ``problemas aritméticos,
geométricos e matriciais'' \emph{e} lógica formal com peso próprio.
\item \textbf{Inglês vale 12 questões, o mesmo que Português.} Interpretação
de texto foi o assunto mais cobrado da prova inteira, nas duas línguas. Para
quem lê documentação técnica em inglês, é a disciplina de melhor retorno por
hora estudada de toda a prova.
\item \textbf{A nota de corte real vai ser alta.} Desenvolvimento de
Software teve 6.257 inscritos em 2024, o perfil mais concorrido de 23.423
candidatos. Os 57,5 pontos mínimos do edital são piso de eliminação, não
meta. A meta é pontuar alto.
\end{enumerate}

\section{Os dois critérios de aprovação --- e o segundo elimina estratégia}

O item \textbf{9.17} do edital exige as duas coisas, \textbf{cumulativamente}:

\begin{enumerate}
\item obter, no mínimo, \textbf{57,5 pontos} na prova objetiva; \emph{e}
\item \textbf{não zerar nenhuma disciplina.}
\end{enumerate}

O segundo critério é o que costuma passar batido, e ele fecha a porta para a
aritmética esperta. Não existe ``abandonar RLM porque são só 5 questões'' nem
``deixar Inglês de lado e compensar nos específicos'': \textbf{um zero em
qualquer uma das seis disciplinas elimina}, com qualquer nota total. Na
prática isso significa garantir pelo menos um acerto seguro em cada bloco de
Módulo I antes de brigar por ponto marginal nos específicos --- e, no dia,
\textbf{não deixar disciplina em branco}, mesmo que seja para chutar as
últimas.

\section{Onde investir o tempo de revisão}

\begin{center}
\begin{tabular}{@{}p{4.6cm}p{2cm}p{6.4cm}@{}}
\toprule
\textbf{Bloco} & \textbf{Peso} & \textbf{Por quê} \\
\midrule
Engenharia de Software & muito alto & Maior bloco em 2024. Requisitos, ágil, testes, ponto de função caem sempre. \\
Banco de Dados & alto & Eixo duplo com Eng.\ Software. No MPU superou tudo. \\
Arquitetura de Software & alto & Microsserviços, REST$\times$SOAP, nuvem, containers, hexagonal. \\
Segurança & alto & ISO 27001/27002:2022, OWASP, OAuth2/SSO, SAST/DAST. \\
Programação & médio-alto & Spring, XML/JSON/REST, DevOps, Git, mobile, low-code. \\
BI & médio & DW, OLAP, ETL, data mining. \\
Governança & médio & ITIL 4, COBIT 2019, Scrum/Kanban, BPMN. \\
Frontend & médio & SPA$\times$PWA, React/Angular/Vue, HTTPS/TLS. \\
Java & médio & Linguagem-base do edital, mas caiu pouco na amostra. \\
Redes & (fora do edital) & Cai mesmo assim; ${\sim}$7,5 pontos históricos. \\
\bottomrule
\end{tabular}
\end{center}

\section{Checklist do dia da prova}

\begin{itemize}
\item Chegar com 1h30 de antecedência. Portões fecham às 12h30, sem
exceção.
\item Documento de identidade original com foto. CPF e certidão de
nascimento \textbf{não} são aceitos.
\item Comprovante de inscrição ou de pagamento.
\item Caneta esferográfica azul ou preta, corpo transparente.
\item \textbf{Proibido:} relógio de qualquer tipo, celular, lápis,
lapiseira, borracha, corretivo, boné, óculos escuros.
\item Lanche e bebida só em embalagem transparente e sem rótulo.
\item Permanência mínima obrigatória: 2 horas.
\item Caderno de questões só pode ser levado se você sair nos últimos 30
minutos.
\item Haverá detector de metais na entrada da sala e dos sanitários.
\end{itemize}
